\section{Tecnología, gobierno y reglas: por qué los ciclos de retroalimentación siempre llegan tarde}

Las empresas tecnológicas convierten en sistemas de software funciones sociales como los pagos, las comunicaciones, la identidad, el transporte, la salud y la IA, y con ello acumulan un poder que antes ejercían instituciones públicas o infraestructuras reguladas. El gobierno debe responder a los daños, pero a menudo va rezagado en talento técnico, visibilidad de los datos y velocidad para iterar sus reglas. Si las empresas ven la regulación como una fricción puramente externa, pierden la oportunidad de dar forma a reglas que puedan cumplirse; si el gobierno solo aplica la ley después de los incidentes, también genera ambigüedad y un efecto inhibidor.

\subsection{Public-private feedback loop}

El ciclo de política pública ideal no es «las empresas innovan primero y el gobierno sanciona al final», sino uno en el que las señales tecnológicas, los daños reales, el diseño de reglas, la capacidad de implementación y la retroalimentación de la industria se corrigen mutuamente de forma continua. El problema de ingeniería central aquí es la observabilidad: el gobierno no ve los riesgos internos de las plataformas y las empresas no conocen las restricciones de la política pública, de modo que cada parte diseña controles con información incompleta.

\lecturefigure{02-public-private-loop.png}{La tecnología y el gobierno necesitan retroalimentación continua, no una rendición de cuentas unidireccional después de los incidentes.}{Subtítulos locales 03:20--06:20; redibujo conceptual.}

\begin{knowledgebox}{Lectura de la figura: la retroalimentación debe volver a la capa de implementación}
Technology produce productos, datos y riesgos nuevos; los observed harms generan señales de los usuarios, de la aplicación de la ley y del mercado; policy plasma los objetivos en leyes, estándares o expectativas regulatorias; implementation los lleva a controls, reporting y feedback. Si el ciclo se detiene en el texto de la política, la organización solo obtiene obligaciones abstractas; si se detiene en la autorregulación empresarial, el público afectado no cuenta con una reparación confiable. Una buena retroalimentación convierte las reglas en controles que pueden probarse y, a la vez, devuelve al diseño de la política los fallos de esos controles.
\end{knowledgebox}

\begin{warningbox}{Tanto «no regular» como «cualquier regulación es buena» son posturas demasiado simples}
Las reglas pueden quedar obsoletas, superponerse o premiar el cumplimiento meramente formal, pero la falta de reglas claras también traslada los costos a los usuarios, a los investigadores y a los equipos de seguridad de primera línea. Lo que el curso pide de verdad es participar en el diseño: definir con claridad a quién se quiere proteger, qué evidencia hay que observar, qué margen de experimentación se permite, cuáles son los umbrales de escalamiento y cuáles los mecanismos de apelación, en lugar de limitarse a expresar apoyo u oposición.
\end{warningbox}

\teachervoice{El ponente enfatiza que, en sus inicios, la industria tecnológica solía repetir como un mantra que «la regulation asfixia la innovación». El giro más maduro que se plantea en clase es que la falta de reglas también genera fricción: las empresas no conocen los límites, los usuarios no tienen confianza y el gobierno solo puede definir los estándares a posteriori, a partir de incidentes y de la aplicación de la ley.}

\subsection{Por qué canales llegan las reglas a las organizaciones}

Las obligaciones de seguridad rara vez provienen de una sola disposición legal. El Congress o las legislaturas estatales pueden crear obligaciones, las agencias administrativas pueden hacer rulemaking o ejercer supervisión, organismos de normalización como el NIST ofrecen marcos operativos, y la aplicación de la ley y los tribunales aplican las reglas generales a hechos concretos. Las empresas multinacionales enfrentan además distintas jurisdicciones, reglas sectoriales, contratos y compromisos con clientes. El sistema de gobernanza debe mantener un obligation inventory, en lugar de que cada incident obligue a buscar sobre la marcha «qué podría aplicar».

\lecturefigure{11-regulation-paths.png}{Las reglas de ciberseguridad se forman por varias vías: legislación, estándares, supervisión, aplicación de la ley y tribunales.}{Subtítulos locales 06:00--09:10; documentos gubernamentales; redibujo conceptual.}

\begin{knowledgebox}{Lectura de la figura: distinguir el origen de la regla y el tipo de evidencia}
Legislation/rules define obligaciones generales; los standards traducen los objetivos a controls y practices; la supervision observa el cumplimiento mediante inspecciones, informes e interacción continua; enforcement/courts emiten juicios sobre expedientes concretos. Los equipos de ingeniería no pueden equiparar un estándar con la ley, ni tomar un solo caso de aplicación de la ley como una norma completa. Lo correcto es que el counsel establezca el juicio de aplicabilidad y que el equipo de seguridad mapee las obligaciones a activos, procesos, registros y evidencia de simulacros.
\end{knowledgebox}

\begin{importantbox}{Obligation-to-control mapping}
Para cada obligación se mantienen cinco columnas: \textbf{source} (origen), \textbf{scope} (qué entidades, datos o sistemas), \textbf{trigger} (cuándo entra en vigor), \textbf{control/evidence} (cómo se demuestra el cumplimiento) y \textbf{owner/reviewer} (quién es responsable y quién lo cuestiona). Así, el incident commander puede ubicar rápidamente los requisitos pertinentes, en lugar de cargarle la investigación jurídica al on-call engineer en un momento de alta presión.
\end{importantbox}

\subsection{Resumen de la sección}

La tensión entre la tecnología y el gobierno es, en esencia, un problema de retroalimentación, capacidad y evidencia. Las organizaciones deben mapear de forma proactiva las leyes, los estándares y la experiencia en aplicación de la ley a controles ejecutables y, al mismo tiempo, hacer que el personal técnico participe en la discusión de políticas, para que las reglas reflejen cómo fallan realmente los sistemas.
